% Chapter 6 - the services.

\guidechapter{THE SERVICES}{%
  \item Files and directories
  \item Loading and saving
  \item The RAM expansion
  \item Processor speed
  \item The screen palette
  \item Network sockets
  \item Fetching over HTTP
}

\section{Files And Directories}

The Ultimate's own filesystem, reached by name. Open a file, read or write it,
seek in it, close it --- the operations you would expect, with two things worth
knowing.

A directory walk is \textbf{one live exchange}: \cmd{ultimate\_opendir()} then
\cmd{ultimate\_readdir()} until it answers \cmd{ULTIMATE\_END}, with no other
command in between. And a filename goes on the wire exactly as you gave it: the
SDK converts nothing, because uppercase PETSCII and uppercase ASCII are the
same bytes and the filesystem does not care about case.

\section{Loading And Saving}

\cmd{ultimate\_load()} takes the firmware's fast path when the machine has one
--- the bytes go straight into C64 memory without passing through the command
interface at all --- and falls back to the ordinary route when it does not. You
do not choose; the SDK does, and tells you where the file ended.

\cmd{ultimate\_bload()} is the one for data: no header consumed, nothing
interpreted, and not one byte past the limit you set.

\section{The RAM Expansion}

Two different things share one name here, and the difference matters.

\textbf{C64 memory to the expansion and back} is the REU's own direct memory
access. There is no command involved: the SDK programs eight registers and the
transfer happens with the processor halted, so by the next instruction it is
already done. Nothing to poll, nothing to wait for.

\textbf{A file to the expansion} is a command, and it is the one that cannot be
done any other way --- the bytes never pass through the C64 at all. Sixteen
megabytes of assets can be loaded without staging them through 64K.

\cmd{ultimate\_reu\_size()} answers how much there is, in 64K banks, and zero
when there is none. It measures rather than asks, because nothing in the
protocol reports it: the expansion aliases, so the first power-of-two boundary
whose write comes back round to the start is the size.

\begin{note}
16MB is the ceiling whatever the machine has, because the expansion's address
registers are 24 bits. A machine with more RAM than that cannot reach the rest
through this interface.
\end{note}

\section{Processor Speed}

Memory-mapped, not a command, and only on an Ultimate~64. The index runs from
0 to 15; 0 is 1MHz and 3 is 4MHz on every machine that has it. Above 3 the same
index means different speeds on different models, so the SDK passes the number
through and does not pretend to know megahertz.

Whether turbo exists at all is the machine owner's decision. Both registers
read \cmd{\$FF} when it has been switched off, which makes availability
testable rather than assumed.

There is a second half to it: bit~7 turns off badlines, which stops the video
chip stealing cycles from the processor. On a fixed loop that is worth about
6.3\% --- measured, and it agrees with the arithmetic.

\section{The Screen Palette}

All sixteen colours, live, on firmware newer than 3.15. This affects the
running palette only --- not flash, not a file --- so a program that crashes
half way through a colour cycle cannot leave the machine looking wrong
permanently.

A whole-palette rotation costs about a quarter of a frame, so colour cycling
directly is comfortably affordable.

\section{Network Sockets}

TCP and UDP, with three things that decide how you use them.

\textbf{A read never waits.} One issued straight after connecting reports
"nothing yet" even when the far end has already replied. Callers poll --- which
is the right answer on a C64, because a blocking read would freeze the machine
for as long as the other end stayed quiet.

\textbf{The end of the stream is \cmd{ULTIMATE\_END}}, the same code a
directory ends with, and it is reported exactly once. After it the handle is
already gone, so do not close it.

\textbf{Connecting can take thirty seconds.} 48 to 75 milliseconds to a host
that is up, and 30.8 seconds to an address with nothing at it before the
firmware gives up.

\begin{listingbox}
for (;;) \{\\
\hspace*{1em}err = ultimate\_net\_read(h, buf, sizeof buf, \&got);\\
\hspace*{1em}if (err == ULTIMATE\_END) break;\hspace{1em}/* far end hung up */\\
\hspace*{1em}if (err != ULTIMATE\_OK)   return err;\\
\hspace*{1em}if (got == 0)             continue;\hspace{1em}/* nothing yet */\\
\hspace*{1em}/* \ldots use buf[0..got-1] \ldots */\\
\}
\end{listingbox}

\section{Fetching Over HTTP}

\cmd{ultimate\_http\_get()} is the whole of the common case: it builds the
request, runs it, and gives the firmware's slot back whether it worked or not.

\begin{listingbox}
uint16\_t got;\\
uint8\_t  err;\\
\\
err = ultimate\_http\_get(url, buf, sizeof buf, \&got);
\end{listingbox}

\cmd{ULTIMATE\_OK} means the server answered below 400. A 404 or a 500 is
\cmd{ULTIMATE\_ERR\_DEVICE} with the number in
\cmd{ultimate\_device\_code()} --- the request itself worked, and the body of
the error page is in your buffer.

\begin{warning}
A URL written as a plain C string literal arrives at the server uppercased.
cc65 translates source characters into PETSCII, and an HTTP path is case
sensitive. Build the URL as bytes, or fold it back at run time.
\end{warning}

\begin{note}
A reply larger than your buffer is truncated and the rest is gone. Nothing in
the protocol offers a range request or a second block, so size the buffer for
what you expect.
\end{note}
